\chapter{Formal Solving Strategies}

\section{Solving a Sudoku Puzzle}

Let us have a Sudoku puzzle with size $n$ and task $T$.
Let $S_{ALL}$ be the union of all valid solutions.

\subsection*{Strategy}

Strategy is an algorithm that takes a state $U_1$ and tries to find a state $U_2$, so that:
$U_1 \succ U_2$ and $U_2 \succeq S_{ALL}$.

If strategy is able to find $U_2$ and explain why all numbers in $U_1 \setminus U_2$
cannot be part of any valid solution, it returns $U_1 \setminus U_2$.
If it is unable to do so, it returns null.

\subsection*{Solving Sequence}

The software is only allowed to progress in solving the Sudoku using available strategies.

The goal of solving is to find a sequence $T = U_0 \succ U_1 \succ \dots \succ U_l \succeq S_{ALL}$,
such that for all $i$, $0 \leq i < l$,
the state $U_{i + 1}$ can be obtained from the state $U_i$ by at least one available strategy, and
it is not possible to obtain any following state from $U_l$ by any available strategy.

\section{Used Strategies}

This section defines all strategies used in the software.

\subsection{Pointing}

Let $R_1$ and $R_2$ be two different Sudoku regions, that have non-empty intersectin.
Since the size of all Sudoku regions is $n$, there are three non-empty parts of $R_1 \cup R_2$ ---
$R_1 \setminus R_2$, $R_2 \setminus R_1$, and $R_1 \cap R_2$.

\begin{thm}[Pointing Strategy]
    If a Sudoku number $m$ cannot be in any cell of $R_1 \setminus R_2$,
    it cannot be in any cell of $R_2 \setminus R_1$.
\end{thm}

\begin{proof}
    According to the Sudoku rules, in each region, the Sudoku number $m$ must appear in exactly one cell.
    Since $m$ cannot be in $R_1 \setminus R_2$, it has to be in $R_1 \cap R_2$,
    hence it cannot be in $R_2 \setminus R_1$.
\end{proof}

\subsection{Perfect Matching Strategies}

This subsection lists all strategies that can be explained by transforming a part of a Sudoku state to
perfect matching of a bipartite graph and using Hall's theorem.

The proof of correctness of these strategy and an algorithm that can find them is in Chapter~\ref{chap:5}.

\begin{defn}
    The \emph{Strategy size} of a Perfect matching strategy is a number $k \in [n - 1]$, where $n$ is the Sudoku size.
\end{defn}

\subsubsection*{Naked Candidates}

If there is $k$ cells in a Sudoku region that together only contain $k$ numbers,
these numbers cannot be in any other cell in the region.

\subsubsection*{Hidden Candidates}

If there are $k$ numbers in a Sudoku region that together only appear in $k$ cells,
these numbers cannot be in any other cell in the region.

\subsubsection*{Generalised X-Wing}

Let $m$ be a Sudoku number and
let $P$ and $Q$ be two sets of $n$ Sudoku regions
such that each cell of the Sudoku appears in exactly one region of each set.
(In the software it is rows+columns, rows+rectangular/jigsaw regions, columns+rectangular/jigsaw regions.)
If all occurrences of the Sudoku number $m$ in $k$ regions from $P$ appear only in $k$ regions in $Q$,
$m$ cannot appear in these regions in $Q$ anywhere else.
(TODO: is this clear?)

\subsection{Edge Label}\label{subsec:formal-edge-label}

This strategy can be used on a Sudoku puzzle with an edge variant.

Let $a$ and $b$ be two edge-adjacent Sudoku cells with a label $l$, $S$ a Sudoku state
and $L_N$ the label-neighbourhood function from the edge variant.
If there exists a number $c \in S_a$ such that no number in $S_b$ can result in a label $l$,
$\forall d \in S_b: (c, d) \notin L_N(l)$,
then the number $c$ can be removed from $S_a$.

% TODO: check again

\subsection{Edge Region}

This strategy can be used on a Sudoku puzzle with an edge variant.

Let $R$ be a Sudoku region, $S$ be a Sudoku state, $E$ be the set of all label-adjacent pairs of cells
and $m \in [n]$ be a Sudoku number.
Let $Q := \{r \in R \mid m \in S_r\}$ be the set of all cells of region $R$ containing the number $m$.
Let $P := \{c \in [n] \times [n] \mid \forall d \in Q:  c = d \lor (c, d) \in E \lor (d, c) \in E\}$ be set of all cells
that are label-adjacent or equal to all cells in $Q$.

In a valid solution, regardless of where in $R$ is the number $m$, any cell $p$ in $P$ must either contain number $m$
or be label-adjacent to a cell containing the number $m$.
Because of that, in a valid solution, the cell $p$ can only contain $m$
or numbers that are allowed by Strategy~\ref{subsec:formal-edge-label} for at least one occurrence of $m$ in $Q$.

% TODO: check again
